\chapter[AI Inputs: English and Visual-State Learning]
{Developing the AI Inputs: English Understanding and Visual-State Learning}
\label{chap:ai-development}

\section{Phase~2 objective and implementation sequence}

Phase~2 extended the Room~315 transport system with two AI-based inputs. Its
objective was to let an operator state a task in English and use the current
paired-camera observation to obtain the rail state needed to plan and complete
the requested shuttle motion. The language path had to produce a task
description compatible with planning, and the image path had to describe the
current rail scene.

Development began with two separate questions: could an English request be
turned into a valid Room~315 task, and could the paired cameras recover the
scene information needed for planning? The two paths were tested separately
before they were connected. The implementation sequence is summarised in
\cref{tab:phase2-sequence}.

\begin{table}[htbp]
\centering
\small
\caption{Phase~2 development sequence.}
\label{tab:phase2-sequence}
\begin{tabularx}{\textwidth}{@{}p{0.08\textwidth}p{0.28\textwidth}X@{}}
\toprule
\textbf{Step} & \textbf{Work performed} & \textbf{Usable result}\\
\midrule
1 & Define language and visual outputs & Semantic proposals contain task fields,
while visual proposals contain scene-state values\\
\tablerowrule
2 & Integrate a local English model and build task validation &
English requests become confirmed \code{TaskGoal} objects\\
\tablerowrule
3 & Generate paired-camera Room~315 data and separate the inputs from labels &
A structured corpus represents shuttle presence, location and payload state\\
\tablerowrule
4 & Partially fine-tune a paired-view ResNet-18 and select a checkpoint &
Two RGB views produce a fixed-order visual-state proposal\\
\tablerowrule
5 & Connect both branches to the runtime checks & A confirmed \code{TaskGoal}
and a current \code{VisualStateObservation} form the
two integration contracts\\
\bottomrule
\end{tabularx}
\end{table}

\section{Reused components and project-specific development}

The two learned components used different adaptation methods. The language
model was integrated as a pretrained, quantised model without fine-tuning on
Room~315 data. The visual network also reused pretrained image features, but it
was subsequently trained on the project's paired-camera corpus.

\begin{table}[htbp]
\centering
\small
\caption{Reused foundations and project-specific Phase~2 work.}
\label{tab:reused-developed-ai}
\begin{tabularx}{\textwidth}{@{}p{0.24\textwidth}p{0.27\textwidth}X@{}}
\toprule
\textbf{Component} & \textbf{Reused foundation} & \textbf{Developed in this
project}\\
\midrule
English understanding & Pretrained Qwen2.5-1.5B-Instruct weights and the local
\code{llama.cpp} inference backend & Saved local configuration, constrained
prompt, strict JSON schema, explicit-fact extraction, field fusion,
clarification, validation and confirmation dialogue\\
\tablerowrule
Visual state & ImageNet-pretrained TorchVision ResNet-18 & Gazebo corpus,
paired-view architecture, 200-value target, masked multi-head loss, partial
fine-tuning, grouped evaluation and integration of the selected checkpoint\\
\tablerowrule
Planning & PlanSys2 and its POPF solver & Room~315 PDDL model, fresh-problem
generation and closed-loop integration (\cref{chap:neurosymbolic})\\
\bottomrule
\end{tabularx}
\end{table}

\section{Building the English task interface}

\subsection{Local model selection and configuration}

The selected semantic model is
\code{Qwen2.5-1.5B-Instruct-GGUF} from the Qwen2.5 family
\citep{qwen25}, using the
\code{qwen2.5-1.5b-instruct-q4\_k\_m.gguf} quantisation. It runs locally on the
CPU through \code{llama.cpp} \citep{llamacpp}; the configured artefact hash is
listed in \cref{app:reproducibility}. The runtime is offline and uses a 4096-token
context, temperature~0, at most 192 generated tokens and seed~315. These
settings record and control the inference conditions and keep each response
short enough for the schema check.

The pretrained weights were used without project-specific language-model
training or fine-tuning. A versioned prompt constrains the general instruction
model to extracting a semantic proposal from an English Room~315 request. The
prompt enumerates the permitted task fields, identities, rails, stations and
slots. It requires one JSON object of type \code{SemanticParseEnvelope} and
excludes PDDL, plan steps, switch updates, stopper updates and shuttle commands.

\subsection{From model output to a confirmed task}

The language path contains four project-specific stages. A rule-based extractor
first records facts stated explicitly, such as \emph{green shuttle},
\emph{right rail} or \emph{slot 2}. Qwen then proposes candidate intent fields
such as transport, inspection, nearest or payload-qualified selection. The
fusion resolver applies fixed precedence, preventing a semantic proposal from
overwriting an explicit user fact or confirmed context. This produces a
\code{TaskGoalDraft}, which is not an executable command. The dialogue manager
requests any missing information, the domain validator checks rail, identity,
station and slot consistency, and the operator confirms the interpretation.
An immutable \code{TaskGoal} is published after these stages are complete.

For example, the request ``Move the green shuttle to slot 2.'' contains a
configured colour alias. The explicit extractor resolves \emph{green} to
\code{room315\_right\_shuttle\_3}; the verb and destination are also grounded
directly from the text. The semantic envelope is stored with the parse record,
and the explicit fields retain precedence. The confirmation therefore states
R3, the right rail and slot~2 before the task is published. The complete
colour-alias execution is followed in \cref{tab:green-shuttle-trace}.

The indirect request ``Could you get a loaded shuttle over to the right
slot~4?'' falls outside the explicit transport-verb patterns. In the saved
evaluation record, the
semantic model supplied \code{goal\_type=transport}, while field fusion and
domain validation constrained the side, payload and destination. A separate
ten-case English corpus covers explicit, indirect, ambiguous, inconsistent and
cancellation forms; its evaluation is reported in
\cref{sec:language-contract-evidence}.

\section{Generating the Room~315 visual corpus}

\subsection{Systematic simulated scenes}

For identities declared present by the configured registry, the visual
branch estimates the location and loaded state required by the planner. The
dataset generator first enumerated every non-empty subset of the eight
canonical identities L1--L4 and R1--R4, giving $2^8-1=255$ presence
configurations. Eight geometrically different scenes were generated for each
configuration by varying valid rail positions, payload states and switch
arrangements, producing 2040 scenarios.

For every scenario, the orchestration layer configured the Gazebo scene,
allowed it to settle and captured one synchronised image from each overhead
rail camera. The resulting corpus contains 2040 paired observations, or 4080
RGB image files. \Cref{fig:paired-camera-training} shows one such model input.

\begin{figure}[htbp]
  \centering
  \begin{subfigure}[t]{0.485\textwidth}
    \includegraphics[width=\textwidth]{f13_left_raw.png}
    \caption{Left-camera RGB input}
  \end{subfigure}\hfill
  \begin{subfigure}[t]{0.485\textwidth}
    \includegraphics[width=\textwidth]{f13_right_raw.png}
    \caption{Right-camera RGB input}
  \end{subfigure}
  \caption[Paired Room~315 camera observation used for visual training.]
  {The two synchronised RGB images form the complete learned input for one
  scenario. Simulator state is used to construct labels in a separate file and
  is never appended to the image tensor.}
  \label{fig:paired-camera-training}
\end{figure}

\subsection{Inputs, labels and grouped partitions}

Gazebo supplies the known identity, calibrated rail coordinate, payload state
and projected image bounding box used to create supervision. Those values are
stored separately from the image references. The model features are the paired
RGB overhead images; depth, controller position, point clouds and row-level
oracle metadata remain label-generation data. This separation
prevents simulator truth from leaking into inference.

All eight variants of one presence configuration remain in the same data
partition. The resulting grouped split contains 191 configurations and 1528
scenarios for training, 32 configurations and 256 scenarios for validation,
and 32 configurations and 256 scenarios for the locked Test partition. Image
hashes and configuration-family membership are audited across the partitions.
Checkpoint selection uses the validation partition alone; the Test partition
is used once for the final evaluation in \cref{chap:results}.

\section{Designing and training the visual-state estimator}

\subsection{Fixed 200-value supervision contract}

Each image pair is resized to $224\times224$ per view and normalised using the
ImageNet channel statistics. The same ResNet-18 backbone \citep{resnet} is
applied to the left and right images. Their two 512-value feature vectors are
concatenated, passed through a 128-unit structured head and mapped to exactly
200 output values. The fixed identity order is L1--L4 followed by R1--R4, with
25 values assigned to each identity:

\begin{equation*}
  25 =
  \underbrace{2}_{\text{rail-side scores}} +
  \underbrace{14}_{\text{block scores}} +
  \underbrace{2}_{\text{payload scores}} +
  \underbrace{4}_{\text{bounding box }(x,y,w,h)} +
  \underbrace{3}_{s_m,\;s_{ratio},\;L_{segment}} .
\end{equation*}

The fixed output index supplies shuttle identity, while a configured registry
specifies which identities exist in the launched scene. At runtime, the model
estimates location and payload values for identities marked present by that
registry.

\begin{figure}[p]
  \centering
  \begin{tikzpicture}[
  font=\sffamily\footnotesize,
  execute at begin node={\hyphenpenalty=10000\exhyphenpenalty=10000},
  input/.style={draw=black!42,rounded corners=2pt,fill=black!3,
    align=center,text width=66mm,minimum height=13mm,inner sep=1.4mm},
  encoder/.style={draw=mfjablue,very thick,rounded corners=2pt,
    fill=mfjalight,align=center,text width=66mm,minimum height=24mm,
    inner sep=1.5mm},
  learned/.style={draw=mfjablue,very thick,rounded corners=2pt,
    fill=mfjablue!5,align=center,text width=152mm,minimum height=14mm,
    inner sep=1.5mm},
  contract/.style={draw=black!55,very thick,rounded corners=2pt,
    fill=white,align=center,text width=152mm,minimum height=30mm,
    inner sep=1.6mm},
  runtimeinput/.style={draw=black!48,rounded corners=2pt,fill=black!3,
    align=center,text width=68mm,minimum height=13mm,inner sep=1.3mm},
  gate/.style={draw=mfjagreen,very thick,rounded corners=2pt,
    fill=mfjagreen!7,align=center,text width=152mm,minimum height=19mm,
    inner sep=1.6mm},
  presentcase/.style={draw=mfjagreen,very thick,rounded corners=2pt,
    fill=mfjagreen!7,align=center,text width=47mm,minimum height=23mm,
    inner sep=1.4mm},
  absentcase/.style={draw=black!58,very thick,rounded corners=2pt,
    fill=black!3,align=center,text width=47mm,minimum height=23mm,
    inner sep=1.4mm},
  rejectcase/.style={draw=mfjared,very thick,rounded corners=2pt,
    fill=mfjared!6,align=center,text width=47mm,minimum height=23mm,
    inner sep=1.4mm},
  accepted/.style={draw=mfjagreen,very thick,rounded corners=2pt,
    fill=mfjagreen!10,align=center,text width=152mm,minimum height=23mm,
    inner sep=1.6mm},
  paneltitle/.style={anchor=west,font=\sffamily\small\bfseries,
    text=mfjablue},
  flow/.style={-{Latex[length=2.1mm]},very thick,draw=mfjablue},
  gateflow/.style={-{Latex[length=2.1mm]},very thick,
    draw=mfjagreen!80!black},
  caseflow/.style={-{Latex[length=1.9mm]},very thick,draw=black!75}
]
  % Panel A: the actual paired-view architecture used in training and runtime.
  \node[paneltitle] at (-79mm,11mm)
    {A \enspace LEARNED PAIRED-VIEW STATE ESTIMATOR};
  \draw[mfjablue!28,line width=0.55pt] (-79mm,7mm) -- (79mm,7mm);

  \node[input] (left) at (-42mm,0) {
    \textbf{Left RGB image}\\
    resize to \(224\!\times\!224\); ImageNet normalisation};
  \node[input] (right) at (42mm,0) {
    \textbf{Right RGB image}\\
    resize to \(224\!\times\!224\); ImageNet normalisation};

  \node[encoder] (leftencoder) at (-42mm,-25mm) {
    \textbf{Shared-weight ResNet-18}\\[-0.2ex]
    left view processed as three-channel RGB\\
    ImageNet-pretrained; early stages frozen\\
    \texttt{layer4} fine-tuned $\rightarrow$ 512-D feature \(f_L\)};
  \node[encoder] (rightencoder) at (42mm,-25mm) {
    \textbf{Shared-weight ResNet-18}\\[-0.2ex]
    right view processed as three-channel RGB\\
    ImageNet-pretrained; early stages frozen\\
    \texttt{layer4} fine-tuned $\rightarrow$ 512-D feature \(f_R\)};

  \node[learned] (concat) at (0,-51mm) {
    \textbf{Feature fusion}\qquad
    \([f_L;f_R]\): \(512+512\rightarrow1024\) values};
  \node[learned,minimum height=17mm] (head) at (0,-69mm) {
    \textbf{Post-fusion prediction head}\\[-0.2ex]
    LayerNorm(1024) $\rightarrow$ Linear(128) $\rightarrow$ SiLU
    $\rightarrow$ Dropout(0.05) $\rightarrow$ Linear(200)};

  \node[contract] (contract) at (0,-95mm) {
    \textbf{Fixed named state proposal: 200 normalised outputs}\\[0.6mm]
    Decoder identities: \texttt{L1 L2 L3 L4 R1 R2 R3 R4}\\[0.4mm]
    Per-identity semantics: \textcolor{mfjablue}{18 categorical scores}
    (side 2, block 14, payload 2)\\[-0.1ex]
    + \textcolor{mfjaorange}{7 numeric targets}
    (bounding box 4, \(s_m\), \(s_{ratio}\), auxiliary \(L_{segment}\))\\[0.3mm]
    \textbf{\(8\times(18+7)=200\)} };

  \draw[flow] (left) -- (leftencoder);
  \draw[flow] (right) -- (rightencoder);
  \draw[flow] (leftencoder.south) -- ([xshift=-42mm]concat.north);
  \draw[flow] (rightencoder.south) -- ([xshift=42mm]concat.north);
  \draw[flow] (concat) -- (head);
  \draw[flow] (head) -- (contract);

  % Panel B: the configured registry decides which identity values may be
  % decoded; the neural network itself never decides presence.
  \node[paneltitle,anchor=center] at (0,-117mm)
    {B \enspace RUNTIME DECODING AND CHECKS};
  \draw[mfjablue!28,line width=0.55pt] (-79mm,-121mm) -- (79mm,-121mm);

  \node[runtimeinput] (proposal) at (-40mm,-133mm) {
    \textbf{Model proposal + selected artefacts}\\
    denormalise with training statistics};
  \node[runtimeinput] (presence) at (40mm,-133mm) {
    \textbf{Current identity registry}\\
    \texttt{ShuttleState} names and timestamps only};

  \node[gate] (decode) at (0,-153mm) {
    \textbf{Identity-aware presence gate}\\[-0.2ex]
    pair each fixed decoder identity with its independent
    \texttt{present}, \texttt{absent} or \texttt{unknown} state};

  \node[presentcase] (present) at (-53mm,-177mm) {
    \textbf{present}\\[-0.2ex]
    decode this identity's learned state values};
  \node[absentcase] (absent) at (0,-177mm) {
    \textbf{absent}\\[-0.2ex]
    ignore its learned values; retain configured absence};
  \node[rejectcase] (reject) at (53mm,-177mm) {
    \textbf{unknown / conflict}\\[-0.2ex]
    reject this observation; publish no visual update};

  \node[accepted] (accepted) at (0,-203mm) {
    \textbf{After all identities are consistent: build the visual observation}\\[-0.2ex]
    check artefacts, freshness, vocabulary and numeric consistency\\
    $\rightarrow$ publish \texttt{VisualStateObservation}};

  % Keep every connector vertical by using aligned ports on the wide nodes.
  % This preserves the two independent inputs without diagonal crossings.
  \draw[flow] ([xshift=-40mm]contract.south) -- (proposal.north);
  \draw[gateflow] (proposal.south) -- ([xshift=-40mm]decode.north);
  \draw[gateflow] (presence.south) -- ([xshift=40mm]decode.north);
  \draw[caseflow] ([xshift=-53mm]decode.south) -- (present.north);
  \draw[caseflow] (decode.south) -- (absent.north);
  \draw[caseflow] ([xshift=53mm]decode.south) -- (reject.north);
  \draw[gateflow] (present.south) -- ([xshift=-53mm]accepted.north);
  \draw[gateflow] (absent.south) -- (accepted.north);

  \path[use as bounding box] (-82mm,14mm) rectangle (82mm,-216mm);
\end{tikzpicture}

  \caption[Paired-view visual model and fixed-order state decoding.]
  {The two RGB views are encoded independently with the same ResNet-18 weights.
  Their 512-value features are concatenated and mapped to a fixed 200-value
  state proposal. A separate configured identity registry controls which
  named values may be decoded; whole-observation validation decides whether a
  \code{VisualStateObservation} is published.}
  \label{fig:visual-model-contract}
\end{figure}

The separate 200-element \code{target\_mask} is used only to compute the loss.
When an identity is absent or lacks valid visual supervision, its 25 target
coordinates receive zero weight; valid present identities contribute normally.
The mask is not used for image segmentation or as a runtime ROS message, and it
remains outside the command path.

\subsection{Training procedure}

Training started from the official
\code{ResNet18\_Weights.IMAGENET1K\_V1} checkpoint. The earlier backbone stages
were frozen, while ResNet \code{layer4} and the custom prediction head were
fine-tuned. This left 8,552,776 of 11,335,560 parameters trainable. Targets were
normalised using statistics computed from valid training coordinates only.

The loss is the equally weighted mean of five masked Smooth-L1 heads:
segment location, loaded state, bounding box, $s_m$ and $s_{ratio}$. Equal
head weighting prevents larger coordinate groups from dominating by element
count, while masking removes absent identities from the loss. The exact
training configuration is summarised in \cref{tab:visual-training-config}.

\begin{table}[htbp]
\centering
\small
\caption{Configuration of the selected visual training run.}
\label{tab:visual-training-config}
\begin{tabularx}{\textwidth}{@{}p{0.31\textwidth}X@{}}
\toprule
\textbf{Item} & \textbf{Configured value}\\
\midrule
Input & Two RGB views, each resized bilinearly to $224\times224$ and
ImageNet-normalised; no data augmentation\\
\tablerowrule
Initial weights and adaptation & ImageNet-pretrained ResNet-18; partial
fine-tuning of \code{layer4} plus the structured head\\
\tablerowrule
Optimiser & AdamW, learning rate $10^{-3}$, weight decay $10^{-4}$\\
\tablerowrule
Batch and schedule & Batch size 32; maximum 15 epochs; validation-loss patience
3\\
\tablerowrule
Loss & Masked per-sample Smooth-L1, equal weight for five output heads\\
\tablerowrule
Reproducibility & CUDA training, seed 31520260730; Test partition excluded from
training and model selection\\
\bottomrule
\end{tabularx}
\end{table}

All 15 epochs completed. Epoch~14 had the minimum total weighted validation
loss and was therefore saved as \code{best.pt}. Its checkpoint, vectoriser,
target statistics, preprocessing description and training configuration were
stored together. At startup, the runtime checks their expected hashes and schema
before the model can publish a proposal. The final locked-Test measurements are
presented in \cref{chap:results}.

\begin{figure}[p]
  \centering
  \resizebox{0.98\textwidth}{!}{\begin{tikzpicture}[
  font=\sffamily\small,
  number/.style={circle,draw=mfjablue,fill=mfjablue,text=white,
    minimum size=8mm,inner sep=0pt,font=\sffamily\small\bfseries},
  title/.style={draw,rounded corners=2pt,align=center,text width=38mm,
    minimum height=10mm,inner sep=1.5mm,font=\sffamily\small\bfseries},
  detail/.style={draw,rounded corners=2pt,align=center,text width=92mm,
    minimum height=10mm,inner sep=1.5mm,font=\sffamily\small},
  teal/.style={draw=mfjagreen,fill=mfjagreen!6},
  blue/.style={draw=mfjablue,fill=mfjalight},
  green/.style={draw=mfjagreen!90!black,fill=mfjagreen!7},
  orange/.style={draw=mfjaorange,fill=mfjaorange!8},
  red/.style={draw=mfjared,fill=mfjared!6},
  flow/.style={-{Latex[length=2mm]},thick,mfjablue}
]
  \node[number] (n1) at (-7.45,0) {1};
  \node[title,teal] at (-4.95,0) {Generate\\varied scenes};
  \node[detail,teal] at (2.25,0)
    {Create eight geometric variants for\\every non-empty identity configuration};

  \node[number] (n2) at (-7.45,-1.45) {2};
  \node[title,teal] at (-4.95,-1.45) {Capture paired images};
  \node[detail,teal] at (2.25,-1.45)
    {Record synchronised left- and right-camera RGB frames for each scene};

  \node[number] (n3) at (-7.45,-3.30) {3};
  \node[title,blue,minimum height=22mm] at (-4.95,-3.30)
    {Separate inputs\\and labels};
  \node[detail,blue,minimum height=22mm] at (2.25,-3.30) {};
  \node[draw=mfjablue,fill=mfjalight,rounded corners=2pt,align=center,
    text width=40mm,minimum height=13mm,font=\sffamily\small]
    at (-0.15,-3.30) {\textbf{Model input}\\Paired RGB only};
  \node[draw=black!50,fill=black!4,rounded corners=2pt,align=center,
    text width=40mm,minimum height=13mm,font=\sffamily\small]
    at (4.65,-3.30) {\textbf{Training labels}\\Simulator state and target mask};

  \node[number] (n4) at (-7.45,-5.15) {4};
  \node[title,green] at (-4.95,-5.15) {Group related scenes};
  \node[detail,green] at (2.25,-5.15)
    {Keep all variants of one identity configuration in the same partition};

  \node[number] (n5) at (-7.45,-7.15) {5};
  \node[title,orange,minimum height=28mm] at (-4.95,-7.15)
    {Train, select\\and evaluate};
  \node[detail,orange,minimum height=28mm] at (2.25,-7.15) {};
  \node[draw=mfjaorange,fill=mfjaorange!7,rounded corners=2pt,align=center,
    text width=39mm,minimum height=11mm,font=\sffamily\small]
    at (-0.20,-6.50) {Train};
  \node[draw=mfjaorange,fill=mfjaorange!7,rounded corners=2pt,align=center,
    text width=39mm,minimum height=11mm,font=\sffamily\small]
    at (4.70,-6.50) {Validation\\model selection};
  \node[draw=mfjared,fill=mfjared!5,rounded corners=2pt,align=center,
    text width=82mm,minimum height=11mm,font=\sffamily\small]
    at (2.25,-7.80) {Final Test evaluation};

  \node[number] (n6) at (-7.45,-9.35) {6};
  \node[title,red,minimum height=18mm] at (-4.95,-9.35)
    {Use the selected\\runtime model};
  \node[detail,red,minimum height=18mm] at (2.25,-9.35)
    {Deploy the selected epoch with its\\preprocessing and output description};

  \draw[flow] (n1) -- (n2);
  \draw[flow] (n2) -- (n3);
  \draw[flow] (n3) -- (n4);
  \draw[flow] (n4) -- (n5);
  \draw[flow] (n5) -- (n6);
\end{tikzpicture}
}
  \caption[Visual-data, training and checkpoint-promotion sequence.]
  {The visual-data workflow from scene generation to the model used at runtime.
  Paired images are kept separate from their training labels, and all variants
  of one identity configuration remain in the same partition.}
  \label{fig:dataset-pipeline}
\end{figure}

\section{Integration outputs}

The resulting integration contracts are a confirmed \code{TaskGoal} and a
current \code{VisualStateObservation}. \Cref{chap:neurosymbolic} follows how
the gateway uses them for planning and shuttle motion.
